\begin{dang}{DAO ĐỘNG TẮT DẦN VÀ HIỆN TƯỢNG CỘNG HƯỞNG}
\begin{itemize}
\item Bài tập về hiện tượng cộng hưởng
\begin{itemize}
\item Hiện tượng cộng hưởng xảy ra khi chu kì dao động cưỡng bức bằng chu kì dao động riêng:  
\begin{align*}
T_{cb}=T_0: \begin{cases}
T_{cb}=\dfrac{\Delta S}{v}=\dfrac{2\pi}{\omega_{cb}}\\
T_0=2\pi \sqrt{\dfrac{m}{k}}=2\pi \sqrt{\dfrac{\ell}{g}}
\end{cases}
\end{align*}
\item So sánh biên độ dao động cưỡng bức:
\begin{itemize}
\item  Xác định vị trí cộng hưởng:
$\omega_0=2\pi f_0=\dfrac{2\pi}{T_0}=\sqrt{\dfrac{k}{m}}$ hoặc $\omega_0=\sqrt{\dfrac{g}{\ell}}$.
\item  Vẽ đường cong biểu diễn sự phụ thuộc biên độ dao động cưỡng bức vào tần số dao động cưỡng bức.
\item  So sánh biên độ và lưu ý: càng gần vị trí cộng hưởng biên độ càng lớn, càng xa vị trí cộng hưởng biên độ càng bé. 
\end{itemize}
\end{itemize}
\end{itemize}

\begin{itemize}
\item Quãng đường vật đi được cho tới khi dừng
\begin{itemize}
\item Cơ năng ban đầu của hệ:$W_0=\dfrac{1}{2}m\omega^2 A^2=\dfrac{1}{2}kA^2$.
\item Vì dao động tắt dần là do cơ năng chuyển hóa thành công của lực ma sát nên 
\begin{align*}
A_{ms}=F_{ms}.s=\mu mgs
\end{align*}
\item Đến khi vật dừng lại thì toàn bộ  \bth{W_0} chuyển hóa thành công của lực ma sát nên 
\begin{align*}
W_0=A_{ms}
\end{align*}
Nên quãng đường vật đi được cho đến khi dừng lại là
\begin{align*}
s=\dfrac{W_0}{\mu mg}=\dfrac{\dfrac{1}{2}\omega^2A^2}{\mu g}=\dfrac{\dfrac{1}{2}kA^2}{\mu mg}.
\end{align*}
\end{itemize}
\end{itemize}
\end{dang}

